\section{The global high-twist bound from integral homology}
\label{sec:sr-global-cohomology}
We prove the quantitative cohomology estimate needed at the exceptional
dyadic twists. The argument uses integral linear-group homology, retaining
the finite stabilizers and hence the small coefficient primes.
The alternative route in Sections 6.76--6.83 of the
\href{https://principia-math.com/api/ws/files/8c527eb7a1dbba0b1dc69c2ba0e60db0}{777-page reorganized working draft}
suggested this approach. We give the geometric count, the integral
torsion calculation and all numerical absorptions explicitly.
The standard comparison theorems used are stated below with their
published sources; no theorem of the Lean archive supplies the estimate.

For a totally imaginary number field \(F\), put
\[
 n=[F:\mathbb Q],\qquad D=|\operatorname{disc}F|,\qquad
 S_\ell=\{v:v\mid\ell\},\qquad
 X_{\ell,a}(F)=H^2(G_{F,S_\ell},\mathbb Z_\ell(a+1)).
\]
Here \(G_{F,S_\ell}\) is the full Galois group unramified outside
\(S_\ell\), and cohomology uses continuous cochains with the compact
coefficient module.

\begin{srtheorem}\label{thm:sr-global-high-twist}
Let \(a\geq1\) and let \(\ell\) be prime.
Suppose \(\mu_\ell\subset F\) if \(\ell\) is odd, and
\(\mu_4\subset F\) if \(\ell=2\).
Then \(X_{\ell,a}(F)\) is finite and
\begin{equation}\label{eq:sr-global-high-twist}
 \log\#X_{\ell,a}(F)
       \leq \exp\bigl((2an)^{24}\log(2D)\bigr).
\end{equation}
The exponent is independent of \(F,\ell,a\).
\end{srtheorem}
The proof occupies the following subsections.

\subsection{The standard comparison theorems}
We use the following established results, in these precise forms.
First, the even group \(K_{2a}(\mathcal O_F)\), \(a\geq1\),
is finite by Quillen's finite generation and Borel's rank theorem.
The arithmetic Quillen--Lichtenbaum comparison gives
\begin{equation}\label{eq:sr-K-comparison}
 K_{2a}(\mathcal O_F)\otimes\mathbb Z_\ell
                      \simeq X_{\ell,a}(F).
\end{equation}
For odd \(\ell\), see
\href{https://www.sas.rochester.edu/mth/sites/doug-ravenel/otherpapers/weibel-ints.pdf}{Weibel, Theorem 70 and Corollary 71}.
For two, the totally imaginary comparison is
\href{https://sites.math.rutgers.edu/~weibel/archive/papers-dir/RognesWeibel.pdf}{Rognes--Weibel, Theorem 0.4}.
Inverting \(\ell\) does not change the \(\ell\)-primary group in
positive even degree: in Quillen's localization sequence the adjacent
positive-degree groups of the finite residue fields have orders prime
to \(\ell\).
Finite-coefficient restricted Galois cohomology agrees with the
corresponding etale cohomology of \(\mathcal O_F[1/\ell]\).
The hypotheses for this comparison are \(S_\ell\) containing all
places above \(\ell\), and no real-place correction at two; see
\href{https://www.mathi.uni-heidelberg.de/~schmidt/papers/kpi1.pdf}{Schmidt, Proposition 2.3 and its proof}.
The finite-coefficient groups are finite, so inverse limits have the
Mittag--Leffler property. The degree-two compact-coefficient comparison
therefore has no \(\lim^1 H^1\) term. These observations specify both
the ramification set and the twist in \eqref{eq:sr-K-comparison}.
In particular this comparison proves finiteness before we estimate
the order.

Second, the plus construction gives
\(K_m(\mathcal O_F)=\pi_m(B\mathrm{GL}(\mathcal O_F)^+)\) for \(m>0\).
The indicated connected component is an infinite loop space, and the
plus map preserves integral homology.
Integral homology stability for number rings identifies its \(m\)-th
homology with
\(H_m(\mathrm{GL}_N(\mathcal O_F),\mathbb Z)\) when
\(N\geq2m+1\); we use the more generous choice
\(N=2m+4\).
This is the number-ring stability theorem of van der Kallen, also used
in
\href{https://www.maths.tcd.ie/EMIS/journals/DMJDMV/vol-kato/soule.dm.pdf}{Soule, Section 4.2}.

Third,
\href{https://msp.org/agt/2016/16-2/agt-v16-n2-p12-p.pdf}{Mathew, Theorem 6.3}
bounds the kernel of the unstable Hurewicz map of a connected infinite
loop space. In degree \(2a\), its \(\ell\)-primary kernel is killed by
\(\ell^{b_\ell}\), where one may take
\[
 b_2=a+3,\qquad
 b_\ell=\left\lceil\frac{2a+2}{2\ell-2}\right\rceil+1
                                             \quad(\ell>2).
\]
Both are at most \(a+3\leq4a\).
The connectivity parameter in that theorem is one; no simple
connectivity of \(B\mathrm{GL}^+\) is required.
We will also use Minkowski's convex body and ideal-class theorems,
the Kummer exact sequence, and the local-invariant form of
Brauer--Hasse--Noether. Their applications and estimates are printed
below.

\subsection{A positive-definite region of forms}
Fix \(N\geq1\), set \(d=nN\), and choose one embedding from each
complex pair. In
\[
 V=\prod_{\sigma:F\hookrightarrow\mathbb C/\mathrm{conj}}\mathbb C^N
                       \simeq\mathbb R^d
\]
use the Euclidean norm
\(\|x\|^2=\sum_\sigma\sum_j|x_{\sigma,j}|^2\), and let
\(\Lambda=\mathcal O_F^N\) with its diagonal embedding.
It is a full lattice with
\[
           \operatorname{covol}(\Lambda)=2^{-d/2}D^{N/2}.
\]
Indeed the determinant of the embedding matrix of an integral basis
squares to the discriminant; replacing each pair of conjugate
coordinates by its real and imaginary parts contributes a factor
\(1/2\) to the absolute determinant.

Every nonzero element of \(\Lambda\) has norm at least one.
A nonzero coordinate \(z\in\mathcal O_F\) satisfies
\(\prod_\sigma|\sigma z|^2=|\operatorname{Norm}_{F/\mathbb Q}z|\geq1\),
so at least one of its archimedean absolute values is at least one.
Let
\[
 W=\prod_\sigma\operatorname{Herm}_N(\mathbb C),\qquad
 q_h(x)=\sum_\sigma x_\sigma^*h_\sigma x_\sigma,\qquad
 P=\{h\in W:q_h(y)\geq1\ (0\ne y\in\Lambda)\}.
\]
The set \(P\) is closed, convex and nonempty: the tuple of identity
matrices belongs to it. Let \(T\subset W\) be the closed cone of
positive semidefinite tuples, and put
\[
 I_0=D^N d^d .
\]

\begin{srlemma}\label{lem:sr-form-region}
Every \(h\in P\) is positive definite and
\begin{equation}\label{eq:sr-form-det}
                       \det_{\mathbb R}h\geq I_0^{-2}.
\end{equation}
Near each point of \(P\), only finitely many lattice inequalities
are needed to define \(P\). A point is a vertex exactly when its
active functionals \(h\mapsto q_h(y)\) span \(W^*\).
\end{srlemma}
\begin{proof}
The rational span \(\Lambda\otimes\mathbb Q\) is dense in \(V\).
If \(q_h\) were negative somewhere, it would be negative on a rational
lattice direction. Clearing the denominator would contradict
\(q_h(y)\geq1\). Thus \(h\in T\).

For \(\delta>0\), the ellipsoid
\(q_{h+\delta\,1}(x)<1\) contains no nonzero lattice point.
If \(v_d\) is the Euclidean unit-ball volume, Minkowski gives
\[
 \frac{v_d}{\sqrt{\det_{\mathbb R}(h+\delta\,1)}}
                      \leq2^d\operatorname{covol}(\Lambda).
\]
The cube \([-1/\sqrt d,1/\sqrt d]^d\) is contained in the unit
ball, so \(v_d\geq2^dd^{-d/2}\).
Dropping the factor \(2^{-d/2}\leq1\) in the covolume formula,
\[
 \sqrt{\det_{\mathbb R}(h+\delta\,1)}
       \geq D^{-N/2}d^{-d/2}=I_0^{-1/2}\geq I_0^{-1}.
\]
Let \(\delta\downarrow0\). The determinant stays positive, hence
the semidefinite tuple \(h\) is positive definite and
\eqref{eq:sr-form-det} follows.

In a sufficiently small neighborhood of such a tuple, all quadratic
forms are bounded below by a common positive multiple of \(\|x\|^2\).
Only a bounded set of lattice vectors can violate an inequality
there, and that set is finite.
Thus \(P\) is locally cut out by finitely many linear inequalities.
If the active functionals at \(h\) do not span \(W^*\), a nonzero
annihilating direction gives a short segment through \(h\) inside
\(P\); the finitely many inactive nearby inequalities have positive
slack. The point is then not extreme.
If they span, any expression of \(h\) as the midpoint of two points
of \(P\) forces all the active inequalities to be equalities at both
points, and their span forces both points to equal \(h\).
This proves the vertex characterization.
\end{proof}

\subsection{The normal-cone complex is integrally acyclic}
Let \(U\subset W^*\) be the open cone of functionals strictly positive
on \(T\setminus\{0\}\).
For each vertex \(v\) of \(P\), set
\[
 C_v=\{u\in U:u(h)\geq u(v)\text{ for all }h\in P\}.
\]
These are convex, relatively closed subsets of \(U\).

\begin{srlemma}\label{lem:sr-normal-acyclic}
The \(C_v\) cover \(U\).
The augmented chain complex freely generated in degree \(p\) by all
ordered tuples \((v_0,\ldots,v_p)\) with
\(C_{v_0}\cap\cdots\cap C_{v_p}\ne\varnothing\), repetitions allowed,
is exact. Its boundary is the alternating sum deleting one entry.
\end{srlemma}
\begin{proof}
For \(u\in U\), compactness of the unit sphere in \(T\) gives
\(u(h)\geq c\|h\|\) on \(T\) for some \(c>0\).
Consequently a sublevel of \(u\) on \(P\) is compact, so its minimum
is attained. The compact minimizing face has an extreme point:
successively minimize independent linear coordinates on that face
until the remaining face is a point. An extreme point of a
minimizing face is extreme in \(P\), hence is a vertex by
Lemma~\ref{lem:sr-form-region}. This proves the cover.

Only finitely many \(C_v\) meet a fixed compact \(K\subset U\).
Fix \(h_0\in P\). On \(K\) the constants \(c\) above have a common
positive lower bound and \(u(h_0)\) has a finite upper bound.
If \(u\in C_v\cap K\), then \(u(v)\leq u(h_0)\); all such \(v\)
therefore lie in a fixed compact subset of \(P\).
Local finite presentation, followed by a finite covering of that
compact set, gives a finite list of possible active inequalities.
Each vertex is specified by a spanning sublist, so only finitely
many occur.

Consider a finite cycle in the ordinary oriented nerve of the
\(C_v\). Choose one witness in the common intersection for every
simplex in its support, and let \(K\) be the convex hull of those
witnesses. This is a compact convex subset of \(U\).
The nonempty \(C_v\cap K\) form a finite closed convex cover of \(K\),
and its nerve contains the cycle.

Here is the finite closed-convex nerve argument used at this point.
For a finite family of compact sets, every subfamily with empty
intersection still has empty intersection after sufficiently small
open thickening. Otherwise a sequence of almost common points would
have a convergent subsequence with limit in the common intersection.
Choose one thickness working for all subfamilies.
The thickenings are open and convex, their intersection pattern is
unchanged, and their union is the corresponding open neighborhood
of \(K\), which is convex. The open convex nerve is acyclic:
the double complex of singular chains on intersections computes
both the nerve and the union, since the intersections are
contractible. Exactness in the covering direction follows after
subdividing each finite singular chain until its simplices lie in
members of the open cover. Thus the original finite cycle bounds.
Every finite cycle in the full ordinary nerve consequently bounds.

Finally compare ordered tuples with the ordinary oriented nerve.
Send a tuple with repeated entries to zero and a tuple with distinct
entries to its sorted oriented simplex. This is a chain map:
the two faces that remove a repeated entry cancel, and on distinct
entries the assertion is the usual orientation rule.
Increasing tuples give a section.
The composite is chain homotopic to the identity on the ordered
complex. To construct the homotopy inductively on a basis tuple,
subtract the already defined homotopy on its boundary from the
difference of the two chain maps.
The result is a cycle supported on the same finite vertex set.
That entire tuple subcomplex contracts by prepending one chosen
vertex, because all its vertices have a common witness.
Use this contraction to fill the cycle and define the next homotopy.
This proves the induction in every degree.
Hence the augmented ordered complex has the same, zero, homology
as the ordinary nerve.
\end{proof}

\subsection{A common frame and a uniform orbit count}
The group \(\Gamma=\mathrm{GL}_N(\mathcal O_F)\) acts by change of
variables on \(P\), preserving the preceding complex.

\begin{srlemma}\label{lem:sr-common-frame}
Every tuple in that complex has \(N\) \(F\)-independent lattice
vectors of value one for every member of the tuple.
If they are the columns of \(A\in M_N(\mathcal O_F)\), then
\[
                   1\leq|\operatorname{Norm}\det A|\leq I_0.
\]
\end{srlemma}
\begin{proof}
Let \(u\) be a common minimizing functional and \(h\) the average
of the tuple. Its active vectors are precisely those active for every
member, since all their individual values are at least one.
If their \(F\)-span were proper, choose a nonzero \(F\)-linear
functional annihilating them.
At one complex place its squared absolute value, and zero at the
other places, defines a nonzero element \(b\in T\).
It vanishes on all active vectors.
Local finite presentation gives \(h-\delta b\in P\) for sufficiently
small \(\delta>0\).
But \(u(b)>0\), contradicting minimality of \(h\).
Thus the common active vectors span \(F^N\), and \(N\) independent
ones can be chosen.

Put \(I=|\operatorname{Norm}\det A|\), a positive integer.
For a member \(h_i\) of the tuple, the pulled-back blocks
\(A_\sigma^*h_{i,\sigma}A_\sigma\) have diagonal entries at most
one: for each column, their sum over the places is one.
Hadamard therefore gives
\(\det_{\mathbb R}(A^*h_iA)\leq1\).
Its real determinant is \(I^2\det_{\mathbb R}h_i\).
Apply \eqref{eq:sr-form-det} to get \(I\leq I_0\).
\end{proof}

Define the positive integers
\[
 R=I_0^2N^d,\qquad V_0=(2R+1)^d,\qquad
 B_p=I_0^{d^2+d+1}(V_0+1)^{dN(p+1)}.
\]
\begin{srproposition}\label{prop:sr-orbit-bound}
The number of \(\Gamma\)-orbits of ordered tuples of \(p+1\) vertices
with a common witness is at most \(B_p\).
Equivalently, \(B_p\) bounds the degree-\(p\) simplex orbits.
For \(L_0=\log(2D)\), and \(m\geq1\), one common bound
\(B=B_{m+1}\) in degrees \(0,\ldots,m+1\) satisfies
\begin{equation}\label{eq:sr-orbit-log}
                   \log B\leq22(m+2)N d^4L_0 .
\end{equation}
\end{srproposition}
\begin{proof}
The row module of the matrix \(A\) in the preceding lemma is an
additive sublattice of \(\mathcal O_F^N\) of index \(I\leq I_0\).
Matrices with the same row module differ by left multiplication by
\(\mathrm{GL}_N(\mathcal O_F)\), since their rows are bases of that
same free module.
Thus these row modules distinguish the possible left orbits of frames.

A fixed integral basis identifies the ambient additive group with
\(\mathbb Z^d\). Hermite normal form bounds the number of its
sublattices of index \(j\) by \(j^{d(d+1)/2}\):
there are at most \(j^d\) diagonal choices and at most
\(j^{d(d-1)/2}\) reduced off-diagonal choices.
Summing over \(1\leq j\leq I_0\) gives at most
\(I_0^{d^2+d+1}\) possibilities.
Counting all additive sublattices only enlarges the number of
\(\mathcal O_F\)-row modules.
Choose one frame representative per orbit and move the tuple so
that its common frame is this representative.

For a fixed representative \(A\), put \(h'=A^*hA\).
All its real eigenvalues are at most \(N\), by the diagonal bound
and the trace bound in each positive Hermitian block.
Their product is at least \(I_0^{-2}\), since \(I\geq1\).
Thus every eigenvalue is at least \(I_0^{-2}N^{-(d-1)}\).
If \(y\in\mathcal O_F^N\) is active for one such vertex and
\(z=A^{-1}y\), then
\[
                    \|z\|\leq I_0N^{(d-1)/2}.
\]
Moreover \(Iz\in\mathcal O_F^N\): the finite additive group
\(\mathcal O_F^N/A\mathcal O_F^N\) has order \(I\), hence is killed
by \(I\). Therefore
\(\|Iz\|\leq I_0^2N^{(d-1)/2}\leq R\).

Distinct points of \(\Lambda\) have distance at least one.
Disjoint open balls of radius \(1/2\) about points in its radius-\(R\)
ball fit inside the radius-\(R+1/2\) ball. Comparing Euclidean volumes
bounds their number by \(V_0\).
The injection \(y\mapsto IA^{-1}y\) therefore leaves at most \(V_0\)
possible active vectors.
A vertex is determined by independent active equations numbering
\(\dim_{\mathbb R}W=nN^2/2\leq dN\).
Its number is at most \((V_0+1)^{dN}\), allowing a dummy entry in
an encoding shorter than \(dN\).
Taking \(p+1\) vertices and multiplying by the number of frames proves
the formula for \(B_p\).

For clarity, the logarithmic estimate also has fixed constants.
Use \(L_0\geq1/2\), \(N\leq d\), and
\(\log t\leq t\) for \(t\geq1\). Then
\[
 \log I_0=N\log D+d\log d\leq3d^2L_0,
 \qquad
 \log R\leq2NL_0+3d^2.
\]
As \(R\geq1\), these give
\[
 \log V_0\leq d(2+2NL_0+3d^2)\leq12d^3L_0,
 \qquad
 \log(V_0+1)\leq13d^3L_0.
\]
Consequently
\[
 \log B\leq
       [9+13(m+2)N]d^4L_0
                 \leq22(m+2)Nd^4L_0,
\]
which is \eqref{eq:sr-orbit-log}.
\end{proof}

\subsection{Integral homology with the stabilizers retained}
\begin{srlemma}\label{lem:sr-stabilizer-size}
Every stabilizer of an ordered simplex is finite of order at most
\(J=3^{d^2}\).
\end{srlemma}
\begin{proof}
The stabilizer of a vertex preserves its positive definite real form.
The image of each vector of a fixed integral basis of \(\Lambda\)
therefore lies in a finite set of lattice vectors of fixed norm.
There are finitely many possible matrices, so this stabilizer is finite.
An ordered-tuple stabilizer is a subgroup of the stabilizer of its
first vertex.

A finite subgroup of \(\mathrm{GL}_d(\mathbb Z)\) injects under
reduction modulo three. Indeed, otherwise its congruence kernel
would contain an element \(g\) of prime order \(p\).
Write \(g=1+3^sA\), where \(s\geq1\) and some entry of \(A\) is
not divisible by three.
If \(p\ne3\), expansion of \(g^p\) modulo \(3^{s+1}\) gives
\(pA\equiv0\pmod3\), a contradiction.
If \(p=3\), expansion modulo \(3^{s+2}\) gives the same contradiction:
\[
 (1+3^sA)^3-1
      =3^{s+1}A+3^{2s+1}A^2+3^{3s}A^3,
\]
and the last two terms are divisible by \(3^{s+2}\).
The action on the additive lattice embeds the finite vertex stabilizer
in \(\mathrm{GL}_d(\mathbb Z)\).
Its reduction consequently embeds in
\(\mathrm{GL}_d(\mathbb F_3)\), of order at most \(3^{d^2}\).
\end{proof}

\begin{srlemma}\label{lem:sr-integral-action-bound}
Let a group \(G\) act on an exact augmented complex whose degree-\(p\)
basis consists of ordered simplices with the alternating deletion
boundary. Suppose that in degrees \(0,\ldots,m+1\) there are at most
\(B\) orbits and stabilizers of order at most \(J\), with \(J\geq2\).
Then \(H_m(G,\mathbb Z)\) is finitely generated, its torsion is finite,
and
\begin{equation}\label{eq:sr-integral-action-bound}
 \log\#H_m(G,\mathbb Z)_{\rm tors}
                 \leq B\log(m+2)+mBJ^m\log J .
\end{equation}
\end{srlemma}
\begin{proof}
Apply a free group resolution to the augmented complex.
Exactness of the augmentation identifies the total homology with
\(H_*(G,\mathbb Z)\). Taking group homology first and using the
permutation bases gives the first-quadrant spectral sequence
\[
 E^1_{p,q}=\bigoplus_{[\sigma]\in G\backslash X_p}
                    H_q(G_\sigma,\mathbb Z)
                  \ \Longrightarrow\ H_{p+q}(G,\mathbb Z).
\]
There is no orientation-character twist: fixing an ordered basis
simplex fixes each of its entries.
This description also follows directly by writing its permutation
module as a sum of induced trivial modules and using the free
resolution on each summand.

For a finite group \(H\) of order at most \(J\), the integral
bar group in degree \(q\) is free of rank at most \(J^q\).
Its cycles have at most that rank, so its homology is generated by
at most \(J^q\) elements.
For \(q>0\), restriction and transfer to the trivial subgroup show
that \(|H|\) kills \(H_q(H,\mathbb Z)\).
Thus
\[
                  \#H_q(H,\mathbb Z)\leq J^{J^q}\quad(q>0).
\]
Every positive-row term in total degree \(m\) consequently has order
at most \(J^{BJ^q}\), where \(1\leq q\leq m\).
The same bound applies to its eventual subquotient.

The row \(q=0\) is the orbit chain complex, free of rank at most
\(B\) in each relevant degree.
Each column of its boundary from degree \(m+1\) has Euclidean norm
at most \(m+2\).
If the boundary has rank \(r\leq B\), the torsion order of its
cokernel is the greatest common divisor of its nonzero
\(r\)-by-\(r\) minors, by Smith normal form.
Hadamard bounds each such minor by \((m+2)^r\).
If \(r=0\), the cokernel is free.
The row-zero homology is a subgroup of that cokernel, so its torsion
order is at most \((m+2)^B\).

For \(r\geq2\), the differential has bidegree \((-r,r-1)\).
There is no incoming differential to \((m,0)\), since its source
would have second coordinate \(1-r<0\).
Thus \(E^\infty_{m,0}\) is a subgroup of \(E^2_{m,0}\), and its
torsion has the same upper bound.
This direction avoids any assertion that a quotient of a free
abelian group must be torsion-free.

All the relevant terms are finitely generated.
The abutment has a finite filtration with these subquotients.
For an exact sequence of abelian groups, torsion in the middle
maps into torsion of the last group with kernel contained in
torsion of the first; its order is at most the product of the
two torsion orders.
Apply this successively to the filtration.
The positive-row terms contribute at most
\(B\sum_{q=1}^mJ^q\log J\leq mBJ^m\log J\).
Together with the row-zero bound this proves
\eqref{eq:sr-integral-action-bound}, finiteness of torsion and finite
generation.
\end{proof}

\begin{srproposition}\label{prop:sr-all-prime-homology}
For \(a\geq1\), \(m=2a\), \(N=4a+4\), and \(z=2an\),
\begin{equation}\label{eq:sr-homology-final}
 \log\#H_{2a}(\mathrm{GL}_{4a+4}(\mathcal O_F),\mathbb Z)_{\rm tors}
                           \leq\exp(z^{23}\log(2D)).
\end{equation}
\end{srproposition}
\begin{proof}
Use the ordered normal-cone complex, with
\(B=B_{m+1}\), \(J=3^{d^2}\), \(d=nN\).
It is exact by Lemma~\ref{lem:sr-normal-acyclic}.
Its orbits and stabilizers have the bounds just proved.
Write \(L_0=\log(2D)\geq1/2\).
Since \(\log J\leq2d^2\) and \(m\geq1\),
\[
 \log(m+2)+mJ^m\log J
       \leq5md^2e^{2md^2}\leq e^{7md^2}.
\]
Here \(\log(m+2)\leq3m\), and \(5md^2\leq e^{5md^2}\).
Therefore \eqref{eq:sr-integral-action-bound} and
\eqref{eq:sr-orbit-log} give
\[
 \log\#H_m(\Gamma,\mathbb Z)_{\rm tors}
 \leq
 \exp\bigl([22(m+2)Nd^4+14md^2]L_0\bigr)
 \leq \exp\bigl(36(m+2)Nd^4L_0\bigr).
\]
Finally \(m+2\leq4a\), \(N\leq8a\), \(d\leq8an\), so
\[
 36(m+2)Nd^4
       \leq 4{,}718{,}592\,a^6n^4
       \leq (2an)^{23}.
\]
The last inequality uses \(4{,}718{,}592<2^{23}\) and \(a,n\geq1\).
This proves \eqref{eq:sr-homology-final}, including at the small
torsion primes.
\end{proof}


\subsection{Arithmetic-group cohomology and finite levels}
\label{sec:sr-arithmetic-cohomology}
The same complex gives the coefficient-uniform cohomology input used
in Section~\ref{sec:spectrum-deformation}.
This application is over \(\mathbb C\); the earlier integral argument
was needed to retain torsion at the small primes.

Keep \(F,N,n,d,D,I_0,V_0\) as above, without any roots-of-unity
hypothesis, and set
\[
 s=nN^2,\qquad T_0=(V_0+1)^{dN},\qquad
 C_0=I_0^{d^2+d+1}T_0^{T_0+1}.
\]
These integers depend only on \(F,N\).

\begin{srproposition}\label{prop:sr-arithmetic-cohomology}
For every finite-dimensional complex representation \(V\) of
\(\Gamma_0=\mathrm{GL}_N(\mathcal O_F)\),
\begin{equation}\label{eq:sr-arithmetic-cohomology}
 \sum_{j\geq0}\dim_{\mathbb C}H^j(\Gamma_0,V)
                                      \leq C_0\dim_{\mathbb C}V.
\end{equation}
The cohomology vanishes for \(j\geq T_0\).
If \(H\subset\mathrm{GL}_N(F)\) is commensurable with \(\Gamma_0\),
the corresponding sum for any finite-dimensional \(H\)-module \(V\)
is at most
\[
 C_0[\Gamma_0:H\cap\Gamma_0]\dim_{\mathbb C}V.
\]
\end{srproposition}
\begin{proof}
Use the ordinary oriented nerve in the proof of
Lemma~\ref{lem:sr-normal-acyclic}.
A simplex has a common frame by Lemma~\ref{lem:sr-common-frame}.
For a fixed frame, the proof of Proposition~\ref{prop:sr-orbit-bound}
bounds all possible vertices by \(T_0\).
Its vertices are distinct, so every simplex has at most \(T_0\)
vertices. The augmented oriented complex is therefore exact and
zero in degrees \(p\geq T_0\).
The number of its degree-\(p\) simplex orbits is at most
\[
               I_0^{d^2+d+1}T_0^{p+1}.
\]
Indeed an ordering of its vertices is an ordered simplex in the
previous count, and forgetting the order maps orbit sets onto orbit
sets.

A stabilizer of an unordered simplex permutes its finite vertex set.
The kernel fixes every vertex and is contained in a finite vertex
stabilizer. Thus the full stabilizer is finite.
It can act on the oriented generator by its permutation sign.
Over \(\mathbb C\), invariants of a finite group are exact: average
a lift over the group to lift any invariant vector through a
surjection. All its positive cohomology consequently vanishes.
Applying a group resolution to the exact oriented complex therefore
computes \(H^*(\Gamma_0,V)\) by a cochain complex whose degree-\(p\)
term has dimension at most
\[
           I_0^{d^2+d+1}T_0^{p+1}\dim V.
\]
More explicitly each orbit contributes the stabilizer invariants of
\(V\) tensored with the orientation character; their dimension is
at most \(\dim V\).
There are finitely many orbits in each degree.
This is the cohomological version of the resolution argument in
Lemma~\ref{lem:sr-integral-action-bound}.
It also proves the asserted vanishing.

A cohomology group has dimension at most its cochain group.
Since \(T_0\geq2\),
\[
 \sum_{p=0}^{T_0-1}T_0^{p+1}
       =\frac{T_0^{T_0+1}-T_0}{T_0-1}
       \leq T_0^{T_0+1}.
\]
Summing the preceding cochain dimensions proves
\eqref{eq:sr-arithmetic-cohomology}.

For the last assertion put \(J=H\cap\Gamma_0\).
Restriction from \(H\) to \(J\) is injective in every degree:
the composition of restriction and the transfer back to \(H\)
is multiplication by \([H:J]\), a nonzero complex scalar.
Shapiro's lemma identifies \(H^*(J,V)\) with
\(H^*(\Gamma_0,\operatorname{Coind}_J^{\Gamma_0}V)\).
The latter coefficient module has dimension
\([\Gamma_0:J]\dim V\).
Apply \eqref{eq:sr-arithmetic-cohomology} and the injectivity.
No action of \(H\) on the original integral lattice is being assumed.
\end{proof}

We now give the finite-orbit and level factors rather than leaving
them as an arithmetic-group hypothesis.
Let \(\widehat{\mathcal O}_F\) be the product of the finite local
integer rings, put
\[
 G_f=\mathrm{GL}_N(\mathbb A_{F,f}),\quad
 G_r=\mathrm{GL}_N(F),\quad
 U=\mathrm{GL}_N(\widehat{\mathcal O}_F),\quad
 U(q)=\ker\bigl(U\longrightarrow
        \mathrm{GL}_N(\widehat{\mathcal O}_F/q\widehat{\mathcal O}_F)\bigr)
\]
for positive integers \(q\).
In particular \(U(1)=U\).
Put \(M_0=2^nD\) and
\[
                  C_*=C_0 M_0^{n+1}(9M_0)^s.
\]
For a complex representation \(V\) of \(G_r\), the action on
functions is
\((\gamma f)(x)=\gamma\cdot f(\gamma^{-1}x)\).

\begin{srproposition}\label{prop:sr-level-cohomology}
For every positive integer \(M\) and every finite-dimensional
complex \(G_r\)-module \(V\),
\begin{equation}\label{eq:sr-level-cohomology}
 \sum_{j\geq0}\dim_{\mathbb C}
 H^j\bigl(G_r,\operatorname{Map}(G_f/U(3M),V)\bigr)
                                   \leq C_* M^s\dim_{\mathbb C}V.
\end{equation}
The constant is chosen before \(M\) and \(V\).
\end{srproposition}
\begin{proof}
First identify the \(G_r\)-orbits in \(G_f/U\).
For \(g\in G_f\), the module
\[
                P_g=F^N\cap g\widehat{\mathcal O}_F^N
\]
is a full \(\mathcal O_F\)-lattice and its completion is
\(g\widehat{\mathcal O}_F^N\).
Here is the elementary approximation in that assertion.
Choose an integer \(b>0\) with
\[
 b\widehat{\mathcal O}_F^N\subset
 g\widehat{\mathcal O}_F^N\subset
 b^{-1}\widehat{\mathcal O}_F^N.
\]
Only finitely many primes occur in these inclusions.
The Chinese remainder theorem identifies
\(b^{-1}\mathcal O_F^N/b\mathcal O_F^N\) with the product of its
local quotients at those primes.
Taking the inverse image of the given subgroups in that finite
quotient constructs \(P_g\), proves that it lies between the two
global lattices, and proves the completion assertion.
Conversely the completion of any full lattice is locally free over
each discrete valuation ring, hence has the form
\(g\widehat{\mathcal O}_F^N\).
Right multiplication by \(U\) leaves this lattice unchanged.
Two such lattices are in the same rational matrix orbit exactly
when they are isomorphic as \(\mathcal O_F\)-modules: an isomorphism
extends uniquely to their \(F\)-spans.

For completeness, the classification needed here is already
elementary over a Dedekind domain.
Project a nonzero coordinate of a full lattice onto its image,
a nonzero fractional ideal \(I\).
An invertible ideal is projective.
Indeed choose finitely many \(x_i\in I,y_i\in I^{-1}\) with
\(\sum_i x_iy_i=1\).
Given a surjection onto \(I\) and lifts \(e_i\) of \(x_i\), the
map \(z\mapsto\sum_i(y_i z)e_i\) splits it.
The coordinate projection therefore splits.
Its kernel is a lattice of rank one less, so induction expresses
the original lattice as a direct sum of fractional ideals.
Proposition~\ref{prop:sr-ideal-lattices} converts that sum to
\(\mathcal O_F^{N-1}\oplus I\), where \(I\) is their product.
Two such normalized modules are isomorphic precisely when their
determinant ideals have the same ideal class.
The necessity follows by taking the top exterior power.
For sufficiency, multiply the final coordinate by a generator of
the ratio of the ideals.
Thus the orbit set is the ideal class group.

Choose integral ideal representatives \(I_\lambda\) with
\(s_\lambda=\#(\mathcal O_F/I_\lambda)\leq M_0\), using
Minkowski, and write
\(P_\lambda=\mathcal O_F^{N-1}\oplus I_\lambda\).
The ideal count in Lemma~\ref{lem:sr-global-generator} gives
\(h_F\leq M_0^{n+1}\); that argument uses no roots of unity.
Choose finite adelic representatives \(g_\lambda\) completing
these lattices.
The fiber over an orbit in \(G_f/U\) has at most
\[
 [U:U(3)]
       \leq\#M_N(\mathcal O_F/3\mathcal O_F)=3^s
\]
orbits at level three.
The ring has \(3^n\) elements because \(\mathcal O_F\) is free of
rank \(n\) over \(\mathbb Z\).
The stabilizer in \(G_r\) of any such orbit representative is
\[
 H_\lambda=G_r\cap g_\lambda U(3)g_\lambda^{-1}.
\]
It is independent of the representative inside \(g_\lambda U\),
since \(U(3)\) is normal in \(U\).
It consists of the automorphisms of \(P_\lambda\) acting identically
on \(P_\lambda/3P_\lambda\).
The equality follows by the completion assertion and the same
Chinese remainder theorem.

The group \(\Gamma_0(3s_\lambda)\) is contained in
\(H_\lambda\cap\Gamma_0\).
To verify this, the finite group \(\mathcal O_F/I_\lambda\) has
order \(s_\lambda\), so
\(s_\lambda\mathcal O_F^N\subset P_\lambda\subset\mathcal O_F^N\).
If \(\gamma\equiv1\pmod{3s_\lambda}\), then
\[
       (\gamma-1)P_\lambda\subset
                    3s_\lambda\mathcal O_F^N\subset3P_\lambda.
\]
The inverse has the same congruence, which proves preservation in
both directions as well as the level-three condition.
The automorphism group of \(P_\lambda\) is commensurable with
\(\Gamma_0\), as both lattices lie between rational scalar multiples
of one another. Explicitly its level-\(3s_\lambda\) kernel on
\(P_\lambda\) is contained in \(\Gamma_0\), and has finite index;
the preceding inclusion gives the finite index on the other side.
Therefore
\[
 [\Gamma_0:H_\lambda\cap\Gamma_0]
                             \leq(3s_\lambda)^s\leq(3M_0)^s.
\]
Proposition~\ref{prop:sr-arithmetic-cohomology} bounds the total
cohomology of every finite-dimensional \(H_\lambda\)-module \(W\)
by \(C_0(3M_0)^s\dim W\).
This proves the required coefficient-uniform stabilizer estimate.

Finally consider \(G_f/U(3M)\to G_f/U(3)\).
Every fiber has \(D_M=[U(3):U(3M)]\) elements.
In fact \(D_M\leq M^s\): write \(u=1+3A\) with
\(A\in M_N(\widehat{\mathcal O}_F)\).
Distinct cosets in this quotient have distinct images of \(A\)
modulo \(M\), since equality of these images is equivalent to
\(u^{-1}u'\equiv1\pmod{3M}\).
There are \(M^s\) possible images.

Over a \(G_r\)-orbit at level three, the function module on its
inverse image is coinduced from its stabilizer, with coefficient
module the functions on this \(D_M\)-point fiber with values in
\(V\). This module has dimension \(D_M\dim V\).
This assertion follows by restricting a function to the fiber at
the chosen orbit representative; its values on other fibers are
then the coinduction function on rational translates.
The diagonal action stated before the proposition is exactly the
action transported by this restriction.
Shapiro's lemma applies with no finiteness condition on the index
in \(G_r\).
There are at most \(h_F3^s\) orbits, so summing their bounds gives
\[
 C_0 h_F3^s(3M_0)^sD_M\dim V
                   \leq C_*M^s\dim V.
\]
The product decomposition over the orbit set is finite.
This proves \eqref{eq:sr-level-cohomology}.
\end{proof}

For the rank-three spectrum over the fixed imaginary quadratic
field used in the character estimate, \(s=18\).
Thus this proves the coefficient-uniform arithmetic cohomology and
finite-level inputs in Section~\ref{sec:spectrum-deformation},
with constants independent of the weight and level.
The separate infinitesimal-character criterion, cuspidal cohomology
injection, and the realization and rigidity of deformation points
must still be checked on the representations to which they are applied.
A dimension estimate alone does not provide that realization.

\subsection{The number of cohomology generators}
Put \(X=X_{\ell,a}(F)\), already finite by
\eqref{eq:sr-K-comparison}, and
\(r=\dim_{\mathbb F_\ell}(X/\ell X)\).

\begin{srlemma}\label{lem:sr-global-generator}
With \(L_0=\log(2D)\),
\begin{equation}\label{eq:sr-global-generator}
                         r\log\ell\leq6n^2L_0 .
\end{equation}
\end{srlemma}
\begin{proof}
Reduction of continuous cochains gives an injection
\[
 X/\ell X\ \hookrightarrow\
                  H^2(G_{F,S_\ell},\mathbb F_\ell(a+1)).
\]
One can verify its kernel without any finiteness assertion.
If an integral two-cocycle reduces to a coboundary, choose a
continuous lift of the finite-valued one-cochain, subtract its
coboundary and divide by \(\ell\).
The result is still a continuous two-cocycle, since multiplication
by \(\ell\) is injective on the coefficient lattice.
Thus the original class was in \(\ell X\).

The roots-of-unity hypothesis makes the mod-\(\ell\) cyclotomic
character trivial. The reduced coefficient module is isomorphic
to \(\mu_\ell\).
On \(R=\mathcal O_F[1/\ell]\), the
\href{https://stacks.math.columbia.edu/tag/03PK}{Kummer exact sequence}
therefore gives
\[
 0\longrightarrow\operatorname{Pic}(R)/\ell
 \longrightarrow H^2_{\rm et}(R,\mu_\ell)
 \longrightarrow\operatorname{Br}(R)[\ell]\longrightarrow0.
\]
The Picard group is a quotient of the ideal class group.
Brauer--Hasse--Noether embeds the last group into
\(\bigoplus_{v\mid\ell}\operatorname{Br}(F_v)[\ell]\):
the local invariants vanish outside \(S_\ell\), and all infinite
places are complex.
Each local group has order \(\ell\); there are at most \(n\)
such places. Consequently
\[
                         r\log\ell\leq\log h_F+n\log\ell.
\]

We include a crude sufficient class-number bound.
Minkowski gives an integral representative in every ideal class
with norm at most \(2^nD\).
The number of ideals of norm \(k\) is at most the coefficient
\(d_n(k)\) of \(\zeta(s)^n\).
To check this locally, for primes above a rational prime \(p\)
with residue degrees \(f_i\), send an ideal exponent vector
\((e_i)\), \(\sum f_ie_i=b\), to an \(n\)-tuple of nonnegative
integers by repeating \(e_i\) in \(f_i\) slots and filling the
remaining slots with zeros. This is injective, because
\(\sum f_i\leq n\).
Multiplying over rational primes gives the assertion.
Also \(d_n(k)\leq k^n\), by counting ordered \(n\)-tuples of positive
factors of \(k\).
Summing up to the integer \(2^nD\) proves
\[
 h_F\leq(2^nD)^{n+1},\qquad
 \log h_F\leq(n+1)^2L_0\leq4n^2L_0.
\]
For odd \(\ell\), \(\mathbb Q(\mu_\ell)\subset F\) implies
\(\ell-1\leq n\). For two, \(\ell\leq n+1\) is automatic.
Hence
\(n\log\ell\leq n\log(n+1)\leq n^2\leq2n^2L_0\).
Combining these inequalities proves
\eqref{eq:sr-global-generator}.
\end{proof}

\subsection{Completion of the global bound}
\begin{proof}[Proof of Theorem~\ref{thm:sr-global-high-twist}]
Put \(m=2a,N=4a+4,z=2an,L_0=\log(2D)\).
The comparison \eqref{eq:sr-K-comparison}, followed by Hurewicz
and homology stability, gives a map
\[
 X\longrightarrow
 H_{2a}(\mathrm{GL}_{4a+4}(\mathcal O_F),\mathbb Z)
                                      \otimes\mathbb Z_\ell .
\]
Flatness of \(\mathbb Z_\ell\) over \(\mathbb Z\) preserves the
kernel description. Mathew's bound makes this kernel a subgroup
of \(X[\ell^{b_\ell}]\), with \(b_\ell\leq4a\).
A finite \(\ell\)-group with \(r\) cyclic factors has at most
\(\ell^{br}\) elements killed by \(\ell^b\), as follows immediately
by writing it as a direct sum of cyclic groups.
The kernel therefore has logarithmic order at most
\(4ar\log\ell\).

The image is finite and lies in the torsion of the target.
By Proposition~\ref{prop:sr-all-prime-homology} its logarithmic
order is at most \(\exp(z^{23}L_0)\).
Lemma~\ref{lem:sr-global-generator} consequently gives
\[
 \log\#X\leq24an^2L_0+\exp(z^{23}L_0).
\]
As \(a,n\geq1\), \(24an^2\leq z^5\leq z^{23}\).
For \(t=z^{23}L_0\geq1\),
\[
 \log\#X\leq t+e^t\leq e^{2t}
                    \leq\exp(z^{24}L_0),
\]
using \(z\geq2\). This is \eqref{eq:sr-global-high-twist}.
Every bound precedes the field, the coefficient prime and the twist.
\end{proof}

\begin{srproposition}\label{prop:sr-K2-proved}
The specialized assertion \(K_2(24)\) of
Definition~\ref{def:sr-K2} holds.
The dyadic common-window theorem
\ref{thm:sr-dyadic-window} is consequently unconditional with
\[
 C_K=24,\qquad
 \eta=\epsilon/48,\qquad
 J=2^{26}(19248/\epsilon)^{24}.
\]
\end{srproposition}
\begin{proof}
The field \(L_7\) is totally imaginary, contains \(\mu_4\), and
has degree 1024. Apply Theorem~\ref{thm:sr-global-high-twist}
with \(\ell=2\), \(F=L_7\), and \(2an=2048a\).
This gives exactly the finiteness and inequality in
Definition~\ref{def:sr-K2}.
The remaining conclusions follow by substitution into
Theorem~\ref{thm:sr-dyadic-window}.
\end{proof}

The \(H^2\) group here is also the one in the archive's
\srcid{LemmaK}: its ramification set is precisely the primes above
\(\ell\), and its compact coefficient is the cyclotomic
\((a+1)\)-st twist.
Changing the algebraic closure identifies these Galois groups;
the factorization of the cyclotomic character through the unramified
quotient is unique. The local embedding choices do not enter this
global \(H^2\).
Thus this bound applies uniformly to the allowed localization data.
The character-counting bounds and connected geometric supply remain
separate tasks. This new prose proof has not been separately
formalized in Lean.
